% Источники: exam/materials/преподаватель/2020/23-03-2020-MODEL L5(1).doc; exam/materials/моделирование.pdf, раздел 31; GitHub HumsterProgrammer/iu9-notes/8sem/Моделирование/Моделирование_лекция-07_05-03-26.md.
\section{Подходы к оцениванию параметров зависимости между ненаблюдаемыми одновременно характеристиками объекта. Метод моментов.}

\subsection*{Контекст}

Характеристики $X$ и $Y$ \emph{не наблюдаются одновременно}: для части объектов известно $X$, для другой части — $Y$. Требуется восстановить параметры совместного распределения или зависимости $Y = f(X, \boldsymbol{\theta})$ по выборкам, полученным в разные моменты.

\subsection*{Метод моментов}

\textbf{Идея:} приравнять теоретические моменты распределения (функции параметра $\boldsymbol{\theta}$) к соответствующим выборочным моментам.

Систему уравнений:
\[
  \mu_k(\boldsymbol{\theta}) = m_k, \qquad k = 1, 2, \ldots, p,
\]
где $\mu_k(\boldsymbol{\theta}) = M[X^k]$ — теоретический момент порядка $k$;
$m_k = \frac{1}{n}\sum_{i=1}^n x_i^k$ — выборочный момент.

Решение системы даёт оценку $\hat{\boldsymbol{\theta}}$.

\textbf{Пример (из лекций)} — степенная зависимость $\xi_2 = a\,\xi_1^{b}$.

Логарифмируем:
\[
  \ln \xi_2 = \ln a + b\ln \xi_1.
\]
Вычисляем выборочные средние и дисперсии логарифмов по раздельным выборкам $x_1,\ldots,x_n$ (для $\xi_1$) и $y_1,\ldots,y_m$ (для $\xi_2$):
\[
  \overline{\ln X} = \frac{1}{n}\sum_{i=1}^n \ln x_i, \qquad
  \overline{\ln Y} = \frac{1}{m}\sum_{j=1}^m \ln y_j,
\]
\[
  s^2_{\ln X} = \frac{1}{n-1}\sum_{i=1}^n(\ln x_i - \overline{\ln X})^2, \qquad
  s^2_{\ln Y} = \frac{1}{m-1}\sum_{j=1}^m(\ln y_j - \overline{\ln Y})^2.
\]
Обозначая
\[
  A=\overline{\ln Y}, \qquad B=\overline{\ln X}, \qquad
  C=s^2_{\ln Y}, \qquad D=s^2_{\ln X},
\]
получаем систему
\[
  A = \ln a + bB, \qquad C=b^2D.
\]
Из нее:
\[
  \hat b = \sqrt{\frac{C}{D}}, \qquad
  \hat a = \exp\!\bigl(A-\hat bB\bigr).
\]

\subsection*{Свойства}

\begin{itemize}
  \item Метод моментов даёт состоятельные оценки.
  \item Метод моментов недостаточно эффективен.
\end{itemize}

\clearpage
